%% kapa: exam=1 round=25 year=2014 period=1 form=A subject=civil_law subject_ko=민법 seq=33 number=33
%% source: Q-Net 감정평가사 기출문제 (cst00309) artlSeq=5101207 file=extracted/1008/5101207/2014년 제25회 감정평가사 1차시험 기출문제/1교시(A형).hwp
%% conversion: kordoc 4.24.51 -> tools/gen1.py (auto); review=proofread-sonnet; flags=none
\begin{question}{33}
관습상 법정지상권에 관한 설명으로 옳지 않은 것은? (다툼이 있으면 판례에 의함)
\begin{choices}
  \item 환지처분으로 인하여 토지와 그 지상건물의 소유자가 달라진 경우, 법정지상권이 그 지상건물의 토지 위에 성립한다.
  \item 법정지상권이 성립한 후 건물을 증ㆍ개축한 경우에도 법정지상권은 존속한다.
  \item 공유토지 위에 건물을 소유하고 있는 토지공유자 중 1인이 자기의 토지 지분만을 매도한 경우, 토지 전체에 관해 법정지상권은 성립할 수 없다.
  \item 법정지상권의 지료를 정하는 때에 법정지상권이 설정된 토지 위에 그 건물이 건립되어 있음으로 인하여 토지소유권이 제한받는 사정은 참작ㆍ평가해서는 안 된다.
  \item 시공회사가 자기 소유의 토지 위에 수위실을 건축한 후 그 부지를 타인에게 매도한 경우, 특별한 사정이 없는 한, 법정지상권이 성립될 수 있다.
\end{choices}
\end{question}
\answer{1}
